\documentclass[10pt,a4paper]{amsart}
\usepackage[margin=19mm]{geometry}
\usepackage{amsmath,amssymb,mathtools}
\usepackage[scr=rsfs,cal=euler]{mathalfa}
\usepackage{xfrac}
\usepackage[unicode,hidelinks,bookmarks=false]{hyperref}
\newcommand{\ie}{i.e.\ }
\newcommand{\cf}{cf.\ }
\newcommand{\resp}{respectively\ }
\newcommand{\Subsection}{\subsection}
\providecommand{\nobreakdash}{\nobreak\hbox{-}}
\setlength{\emergencystretch}{2em}
\multlinegap0pt
\usepackage{bm}
\usepackage{bbm}
\usepackage{dsfont}
\usepackage{xspace}
\usepackage{cancel}
\usepackage{mathtools}
\usepackage{amsthm}
\makeatletter
\@ifundefined{proposition}{\newtheorem{proposition}{Proposition}}{}
\@ifundefined{definition}{\newtheorem{definition}[proposition]{Definition}}{}
\@ifundefined{lemma}{\newtheorem{lemma}[proposition]{Lemma}}{}
\@ifundefined{theorem}{\newtheorem{theorem}[proposition]{Theorem}}{}
\makeatother
\newcommand{\N}{\mathbb N}
\newcommand{\R}{\mathbb R}
\newcommand{\mf}[1]{\mathfrak{#1}}
\newcommand{\rgeq}{\trianglerighteq}
\title[Barycentric decompositions for extensive monotone divergence]{Barycentric decompositions for extensive monotone divergences}
\author{Erkka Haapasalo}
\date{Source: arXiv:2509.18725v1 (CC BY 4.0)}
\begin{document}
\maketitle

\noindent\emph{Source and licence.} Barycentric decompositions for extensive monotone divergences. Version arXiv:2509.18725v1 (CC BY 4.0), distributed under the Creative Commons Attribution 4.0 International licence, as recorded in the arXiv licence metadata for this exact version and shown on the abstract page https://arxiv.org/abs/2509.18725v1. Full rights notice: licenses/arxiv-2509.18725-CC-BY-4.0.txt.\par\smallskip

\noindent\textit{Editorial scope.} Counted unit: the Lemma whose source label is \texttt{lemma:key} in the article (source0.tex lines 405--407, source label \texttt{lemma:key}), delivered together with the article's own complete original proof of it (source0.tex lines 409--423, one \texttt{proof} environment of 15 lines). No mathematics is written, repaired or re-derived: statement and proof are the article's own text line for line. Required context retained uncounted: def:deg (lines 313--319); thm:Vergleichsstellensatz (lines 343--358); eq:FritzConditions (lines 345--347). Every definition, display and label of those lines is retained unchanged. Omitted from this package and not redistributed: 1--312, 320--342, 348--404, 408--408, 424--1360. Nothing inside a retained excerpt is omitted or altered: every retained body is delivered exactly as the article's lines are written, so no omission receipt of any kind accompanies this package. Citations: the retained excerpts carry no citation command at all, so no bibliography entry of the article is transcribed and no reference is redistributed. Reference adaptation: none was needed and none was made; every reference inside the delivered unit resolves inside it, so no unresolved reference or citation is produced. Printed slips kept literal: no printed word, symbol, hypothesis or number of the counted statement or of its proof was altered. Layout and class: the article's own class is \texttt{amsart} with packages including geometry, amsaddr, amsthm, amssymb, amsmath, mathtools, amsfonts, bbm, xcolor, sidecap, enumitem, mathabx, tikz-cd, caption; the wrapper is the lane's pinned \texttt{amsart} editorial wrapper at 10pt on A4, which restates the theorem-like environments the retained text needs and adds the article's own macro definitions that the retained text needs (\texttt{\textbackslash N}, \texttt{\textbackslash R}, \texttt{\textbackslash mf}, \texttt{\textbackslash rgeq}), with extra wrapper packages (bm, bbm, dsfont, xspace, cancel, mathtools). The delivered numbering is the unit's own. Assets: no retained span contains a figure environment, a graphics-inclusion command, a PDF-page inclusion, a table or a TikZ picture; no image file is copied into this package, the package declares no asset dependency and no image file is redistributed with it. Licence: version arXiv:2509.18725v1 of this article is distributed under the Creative Commons Attribution 4.0 International licence, as recorded in the arXiv licence metadata for this exact version and shown on the abstract page https://arxiv.org/abs/2509.18725v1. A full rights notice is delivered at licenses/arxiv-2509.18725-CC-BY-4.0.txt. Characters: every retained excerpt is pure ASCII, so the delivered engine, XeLaTeX with its default Latin Modern text font, reports no missing character.
\begin{definition}\label{def:deg}
We say that a preordered semiring $S$ is {\it of degeneracy $d$} for some $d\in\N$ if there is a surjective homomorphism $\|\cdot\|:S\to\R_{>0}^d\cup\{(0,\ldots,0)\}$ with trivial kernel such that
$$
a\rgeq b\ \Rightarrow\ \|a\|=\|b\|,\quad \|a\|=\|b\|\ \Rightarrow\ a\sim b.
$$
In this situation, we denote the component homomorphisms of $\|\cdot\|$ by $\|\cdot\|_{(k)}$ for $k=1,\ldots,d$.
\end{definition}

\begin{theorem}\label{thm:Vergleichsstellensatz}
Let $S$ be a zerosumfree preordered semidomain of polynomial growth and of degeneracy $d$ for some $d\in\N$. We fix a power universal $u\in S$ and the vector-valued homomorphism $\|\cdot\|$ of Definition \ref{def:deg} to define the test spectrum $\hat{\mf D}(S)$. Suppose that $x,y\in S\setminus\{0\}$ are such that $\|x\|=\|y\|$. If
\begin{equation}\label{eq:FritzConditions}
\Delta(x)>\Delta(y)\qquad\forall\Delta\in\hat{\mf D}(S),
\end{equation}
then
\begin{itemize}
\item[(a)] $x^nu^k\rgeq y^nu^k$ for some $k\in\N$ when $n\in\N$ is sufficiently large and
\item[(b)] $xz\rgeq yz$ for some $z\in S\setminus\{0\}$ which can be chosen according to
$$
z=u^k\sum_{\ell=0}^n x^\ell y^{n-\ell}
$$
for some $k\in\N$ and $n\in\N$ sufficiently large.
\end{itemize}
If $x$ is a power universal, we may omit the appearance of $u^k$ in items (a) and (b) above. Conversely, if condition (a) or (b) above holds, then the inequalities in \eqref{eq:FritzConditions} hold non-strictly.
\end{theorem}

\begin{lemma}\label{lemma:key}
Let $D:S\setminus\{0\}\to\R_+$ be a monotone divergence. If $x,y\in S\setminus\{0\}$, $x\sim y$, are such that $\Delta(x)\geq\Delta(y)$ for all $\Delta\in\hat{\mf D}$, then also $D(x)\geq D(y)$.
\end{lemma}

\begin{proof}
Suppose first that $x,y\in S\setminus\{0\}$, $x\sim y$, are such that $\Delta(x)>\Delta(y)$ for all $\Delta\in\hat{\mf D}$. According to Theorem \ref{thm:Vergleichsstellensatz}, this means that there are $n,k\in\N$ such that $u^k x^n\rgeq u^k y^n$. Using the monotonicity of $D$, we now have
$$
kD(u)+nD(x)=D(u^k x^n)\geq D(u^k y^n)=kD(u)+nD(y)\ \Rightarrow\ D(x)\geq D(y).
$$
Assume next that $x,y\in S\setminus\{0\}$, $x\sim y$, are such that $\Delta(x)\geq\Delta(y)$ for all $\Delta\in\hat{\mf D}$. Let us make the inequality strict by adding the power universal. Indeed, for all $m\in\N$, we have
$$
\Delta(ux^m)=\Delta(u)+m\Delta(x)=1+m\Delta(x)>m\Delta(x)\geq m\Delta(y)=\Delta(y^m)
$$
for all $\Delta\in\hat{\mf D}$. According to what we proved earlier, this means that $D(ux^m)\geq D(y^m)$, i.e.,
$$
mD(x)+D(u)\geq mD(y)\ \Rightarrow\ D(x)+\frac{1}{m}D(u)\geq D(y).
$$
Since this holds for all $m\in\N$, we have $D(x)\geq D(y)$.
\end{proof}

\end{document}
